\section{Security Capability Matrix}
\label{sec:capability-matrix}

The threat model is easiest to navigate when attacker access and information
object are separated. Rows in \cref{tab:capability-matrix} describe what part of
the stack the attacker can influence. Columns describe what the attacker is
trying to obtain or preserve. Every cell is a hyperlink to a compact evidence
record in \cref{app:capability-registry}.

\begin{claimbox}[Intuitive rule motivating the matrix]
\emph{``You have to start thinking about everything as a system instead of
individuals.''} A controlled gate can resemble an ``if this, then that,'' but
when the control is in superposition the condition is applied coherently across
the joint state. The relevant security question is therefore not merely which
single wire changed, but where recoverable information resides before and after
the complete transformation.
\end{claimbox}

\subsection{Legend and navigation}
\hypertarget{matrix:legend}{}

\noindent
\riskbadge{RiskGreen}{F}\ Demonstrated or established feasible\quad
\riskbadge{RiskAmber}{P}\ Partial, probabilistic, access-dependent, or disturbance-limited\quad
\riskbadge{RiskBlue}{H}\ Technically plausible but unvalidated here\quad
\riskbadge{RiskRed}{N}\ Exact objective blocked under the stated assumptions\quad
\riskbadge{RiskGray}{X}\ Unnecessary, superseded, or not applicable.

\smallskip
\noindent Symbols add meaning independently of color:
\ourresult\ contribution demonstrated in this work;
\candidatepriority\ candidate literature-priority claim requiring a complete
related-work review;
\workaround\ an alternative route avoids the direct limitation;
\strongaccess\ stronger or combined attacker access is required;
\nogo\ formal no-go boundary; and
\mitigation\ an established or straightforward mitigation exists.
A star therefore marks a project contribution, not an automatic claim of
world-first novelty.

\begin{landscape}
\begin{table}[p]
\centering
\hypertarget{matrix:overview}{}
\caption{Hyperlinked threat--information capability matrix. Select a cell code
for its complete record in \cref{app:capability-registry}. Short phrases state
the direct objective; arrows in the symbol line indicate that a nearby objective
or stronger access model may avoid the direct limitation.}
\label{tab:capability-matrix}
\setlength{\tabcolsep}{2.3pt}
\renewcommand{\arraystretch}{1.02}
\begin{tabular}{>{\raggedright\arraybackslash}p{0.105\linewidth}
                *{6}{>{\raggedright\arraybackslash}p{0.1415\linewidth}}}
\toprule
\cellcolor{DeepBlue!10}\textbf{Access} &
\cellcolor{DeepBlue!10}\textbf{1. Orthogonal / classical value} &
\cellcolor{DeepBlue!10}\textbf{2. Result / syndrome} &
\cellcolor{DeepBlue!10}\textbf{3. Preparation label} &
\cellcolor{DeepBlue!10}\textbf{4. Arbitrary unknown state} &
\cellcolor{DeepBlue!10}\textbf{5. Metadata / ciphertext} &
\cellcolor{DeepBlue!10}\textbf{6. Delayed-use correlation} \\
\midrule
\cellcolor{DeepBlue!6}\textbf{A}\\Channel only &
\matrixsummary{A1}{RiskGreenSoft}{\riskbadge{RiskGreen}{F}}{Read or copy in a known orthogonal basis.}{\mitigation} &
\matrixsummary{A2}{RiskAmberSoft}{\riskbadge{RiskAmber}{P}}{Obtain only if the result is exposed onto the channel.}{\workaround\strongaccess\mitigation} &
\matrixsummary{A3}{RiskAmberSoft}{\riskbadge{RiskAmber}{P}}{Discriminate a restricted known ensemble with errors or inconclusive runs.}{\workaround} &
\matrixsummary{A4}{RiskRedSoft}{\riskbadge{RiskRed}{N}}{Perfect extraction with perfect preservation is blocked.}{\nogo\workaround} &
\matrixsummary{A5}{RiskGreenSoft}{\riskbadge{RiskGreen}{F}}{Copy exposed metadata or ciphertext; plaintext need not follow.}{\mitigation} &
\matrixsummary{A6}{RiskAmberSoft}{\riskbadge{RiskAmber}{P}}{Store a noisy result now and interpret it after disclosure.}{\workaround} \\
\midrule
\cellcolor{DeepBlue!6}\textbf{B}\\Runtime / routing &
\matrixsummary{B1}{RiskGreenSoft}{\riskbadge{RiskGreen}{F}}{Latch an orthogonal intermediate with controlled gates or routing.}{\mitigation} &
\matrixsummary{B2}{RiskGreenSoft}{\riskbadge{RiskGreen}{F}}{Route a measured branch into reclaimed workspace.}{\ourresult\candidatepriority\strongaccess\mitigation} &
\matrixsummary{B3}{RiskAmberSoft}{\riskbadge{RiskAmber}{P}}{Route a label only when an orthogonal label register already exists.}{\strongaccess} &
\matrixsummary{B4}{RiskAmberSoft}{\riskbadge{RiskAmber}{P}}{Move the state exactly, but do not retain a second perfect copy.}{\workaround\nogo} &
\matrixsummary{B5}{RiskGreenSoft}{\riskbadge{RiskGreen}{F}}{Route flags, tags, ciphertext, or control data into hidden workspace.}{\mitigation} &
\matrixsummary{B6}{RiskGreenSoft}{\riskbadge{RiskGreen}{F}}{Retain a branch whose meaning increases after side information arrives.}{\ourresult\candidatepriority\strongaccess} \\
\midrule
\cellcolor{DeepBlue!6}\textbf{C}\\Compiler / transformation &
\matrixsummary{C1}{RiskGreenSoft}{\riskbadge{RiskGreen}{F}}{Insert a coherent latch for an orthogonal secret-dependent value.}{\mitigation} &
\matrixsummary{C2}{RiskBlueSoft}{\riskbadge{RiskBlue}{H}}{Insert covert measurement, routing, or logging operations.}{\candidatepriority\mitigation} &
\matrixsummary{C3}{RiskGreenSoft}{\riskbadge{RiskGreen}{F}}{Re-prepare a state when symbolic source parameters are visible.}{\strongaccess\workaround} &
\matrixsummary{C4}{RiskRedSoft}{\riskbadge{RiskRed}{N}}{No universal clone; redirect, probe, or restrict the ensemble instead.}{\nogo\workaround} &
\matrixsummary{C5}{RiskGreenSoft}{\riskbadge{RiskGreen}{F}}{Leak circuit shape, scheduling, pulse, or workload metadata.}{\mitigation} &
\matrixsummary{C6}{RiskBlueSoft}{\riskbadge{RiskBlue}{H}}{Preserve a latent correlation for later protocol information.}{\candidatepriority\strongaccess} \\
\midrule
\cellcolor{DeepBlue!6}\textbf{D}\\Source / preparation &
\matrixsummary{D1}{RiskGreenSoft}{\riskbadge{RiskGreen}{F}}{The source already possesses the classical preparation value.}{} &
\matrixsummary{D2}{RiskAmberSoft}{\riskbadge{RiskAmber}{P}}{Bias, tag, or correlate outcomes; a genuinely random future result is not known in advance.}{\workaround} &
\matrixsummary{D3}{RiskGreenSoft}{\riskbadge{RiskGreen}{F}}{Prepare a duplicate from known controls instead of cloning an unknown state.}{\ourresult\candidatepriority\workaround\strongaccess} &
\matrixsummary{D4}{RiskRedSoft}{\riskbadge{RiskRed}{N}}{If the source itself lacks the state description, no-cloning still applies.}{\nogo\workaround} &
\matrixsummary{D5}{RiskGreenSoft}{\riskbadge{RiskGreen}{F}}{Copy basis choices, identifiers, parameters, or ciphertext at the source.}{\mitigation} &
\matrixsummary{D6}{RiskGreenSoft}{\riskbadge{RiskGreen}{F}}{Retain the label or measurement record until basis disclosure.}{\ourresult\strongaccess} \\
\midrule
\cellcolor{DeepBlue!6}\textbf{E}\\Hardware / backend &
\matrixsummary{E1}{RiskAmberSoft}{\riskbadge{RiskAmber}{P}}{Preserve and repeatedly observe a selected QND observable, not the full state.}{\mitigation} &
\matrixsummary{E2}{RiskGreenSoft}{\riskbadge{RiskGreen}{F}}{Log or leak outcomes and syndromes handled by the backend.}{\mitigation} &
\matrixsummary{E3}{RiskAmberSoft}{\riskbadge{RiskAmber}{P}}{Infer encoder settings through probes or implementation side channels.}{\mitigation} &
\matrixsummary{E4}{RiskRedSoft}{\riskbadge{RiskRed}{N}}{Hardware access does not permit perfect unknown-state extraction and preservation.}{\nogo\workaround} &
\matrixsummary{E5}{RiskGreenSoft}{\riskbadge{RiskGreen}{F}}{Exploit timing, pulse behavior, calibration, or structural metadata.}{\mitigation} &
\matrixsummary{E6}{RiskAmberSoft}{\riskbadge{RiskAmber}{P}}{Store classical traces; long-lived hidden quantum correlations face noise and reset.}{\mitigation} \\
\bottomrule
\end{tabular}
\end{table}
\end{landscape}

\subsection{The starred project cluster}
\label{subsec:matrix-starred-cluster}
\hypertarget{matrix:starred}{}

The three most important project-specific cells form a chain rather than three
independent claims:

\begin{tcolorbox}[
  enhanced,colback=SoftTeal,colframe=AdvancedTeal,boxrule=0.55pt,
  arc=1.5pt,left=8pt,right=8pt,top=8pt,bottom=8pt,
  title={Source description $\rightarrow$ retained branch $\rightarrow$ later value}]
\centering
\begin{tikzpicture}[
  node distance=1.05cm,
  every node/.style={font=\footnotesize,align=center},
  box/.style={rounded corners=2pt,draw=AdvancedTeal!80!black,
    fill=white,minimum width=3.85cm,minimum height=1.28cm,inner sep=5pt},
  arr/.style={-{Latex[length=2.2mm]},very thick,draw=AccentBlue}
]
\node[box] (d3) {\textbf{\hyperlink{cell:D3}{D3}}\ourresult\candidatepriority\\Duplicate preparation from\\known source controls};
\node[box,right=of d3] (b2) {\textbf{\hyperlink{cell:B2}{B2}}\ourresult\candidatepriority\\Route the measured branch\\into the fifth wire};
\node[box,right=of b2] (b6) {\textbf{\hyperlink{cell:B6}{B6}}\ourresult\candidatepriority\\Interpret the retained record\\after basis information};
\draw[arr] (d3) -- node[above,font=\scriptsize]{clean victim branch} (b2);
\draw[arr] (b2) -- node[above,font=\scriptsize]{side information} (b6);
\end{tikzpicture}
\end{tcolorbox}

In the reconstructed fixed-input model, source access avoids the impossible task
of cloning an arbitrary unknown channel state. The attacker instead prepares the
victim branch from the same classical controls, routes the measured original
branch to $q_4$, and later interprets that record conditioned on basis
information. The chain combines a source-level alternative route with a
runtime-level retention mechanism. It remains explicitly distinct from a
channel-only attack and from a proven bypass of the original detector's exact
Boolean acceptance register.

\subsection{Touched, untouched, and avoided objectives}

A cell is classified against its \emph{direct} objective. A red cell therefore
does not imply that every nearby attack is impossible. It means that the stated
combination of target, access, and preservation requirement cannot be achieved
exactly. The \workaround\ symbol points reviewers to a changed objective or
assumption, such as accepting disturbance, allowing an inconclusive branch,
using multiple copies, moving rather than copying the state, targeting an
orthogonal observable, or obtaining the classical preparation description
upstream. The detailed registry records four fields for each cell:

\begin{enumerate}[leftmargin=1.8em]
  \item the mechanism or direct objective;
  \item what the mechanism necessarily touches;
  \item what may remain untouched under the stated model; and
  \item the limitation, alternative route, mitigation, and evidence status.
\end{enumerate}

This distinction supports an optimization view of the broader research problem.
For protected variable $X$, hidden register $E$, victim-visible subsystem $V$,
and attacker transformation $\mathcal{A}$, one may ask
\begin{equation}
  \max_{\mathcal{A}}\; I_{\mathrm{acc}}(X:E)
  \quad\text{subject to}\quad
  D\!\left(\rho_V^{\mathcal{A}},\rho_V^{\mathrm{expected}}\right)\leq\varepsilon,
  \qquad
  P_{\mathrm{accept}}^{\mathcal{A}}\geq 1-\delta.
  \label{eq:matrix-security-optimization}
\end{equation}
The matrix supplies a structured domain for that future formalization; it is not
yet the optimization's solution.
